% problem-000201 4 铂硫多核配离子
\section*{4. 铂硫多核配离子}
向 $\mathrm { H _ { 2 } P t ( O H ) } _ { \mathfrak { c } }$ _{6}⽔溶液中加⼊⾜量的 $\mathrm { N a _ { 2 } S _ { 2 } O _ { 3 } }$ 溶液，调整pH⾄11.4，在⾼压⽔热釜中 $1 5 0 ~ ^ { \circ } \mathrm { C }$ 反应 17h，得到⼀种深棕⾊的晶体(A·�H_{2}O)。A是⼀种钠盐，式量为1359.4，A的阴离⼦B是⼀种含三个Pt原⼦的多核络离⼦，仅含Pt、S、O三种元素，B中所有Pt的化学环境⼀致，配位原⼦均为S且Pt的配位数与典型的单核Pt配合物相同，存在两种化学环境不同的S，不存在Pt-Pt键和S-S键。

\subsection*{4-1}
通过计算和分析确定络离⼦B的组成。

\subsection*{4-2}
画出B的结构式。

\subsection*{4-3 写出⽣成配合物A的化学⽅程式。}
